% ============================================================
%  THE ORTYX PROTOCOL
%  Part II — Tension
%  Chapter 8: Definitions That Cannot Be Wrong
% ============================================================

\chapter{Definitions That Cannot Be Wrong}
\label{ch:definitional-closure}

\chapepigraph{The most dangerous sentence in theoretical work is:
``By definition, X is Y.'' Not because definitions are wrong,
but because they are immune to evidence. A framework built
primarily of definitions has acquired invulnerability
at the cost of content.}{Flyxion}

\noindent
A definition cannot be wrong. This is its great strength and
its great danger. Definitions are constitutive: they establish
what a term means within a framework rather than making claims
about the world that might turn out to be false. The statement
``jitter is defined as $|T_i - \bar{T}_i|$'' cannot be
challenged by empirical evidence, because it is not making
an empirical claim. It is establishing a stipulation.

The danger arises when conclusions that should be derived
from empirical claims are instead embedded in definitional
choices. If ``meaning'' is defined as ``temporal coherence,''
then ``temporal coherence is necessary for meaning'' is not
a discovery but a tautology. The sentence looks like an
empirical claim about the relationship between temporal
structure and semantic content; it has the grammatical
form of a synthetic proposition. But its truth is guaranteed
by the definition, not by the world. No observation could
falsify it, because any observation of meaning would, by
definition, be an observation of temporal coherence.

\section{Definitional Closure in the Phoenix Corpus}
\label{sec:closure-phoenix}

The Phoenix corpus exhibits several instances of this pattern,
which range from mild to significant in their epistemic
consequences.

The mildest instance concerns the concept of ``authenticity.''
In Chapter~5, the corpus's use of this term was noted as
doing multiple kinds of work simultaneously: aesthetic,
phenomenological, engineering, and ontological. Examining
the corpus's usage more carefully reveals that ``authenticity''
tends to be applied to signals with high temporal coherence
and withheld from signals with low temporal coherence. This
creates the appearance of a discovered connection: authentic
audio is temporally coherent audio. But if authenticity is
implicitly defined as temporal coherence --- if ``authentic''
means ``temporally organized in the way that the Marine
detects'' --- then the connection is not a discovery but
a consequence of the definition.

A more significant instance concerns the concept of ``meaningful
structure.'' The corpus consistently treats temporally coherent
signals as more meaningful than temporally incoherent signals.
In passages where this treatment is explicit, it appears as
an empirical claim: coherent signals carry more information,
more emotional content, more experiential richness. But in
passages where the treatment is implicit, it functions as a
definition: ``meaningful structure'' just is what the Marine
detects; signals that the Marine rates as low-salience are,
by definition, less meaningfully organized.

\begin{auditprop}
\label{ap:definitional-closure}
Definitional Closure~($\FailGeo{2}$) occurs when a framework
defines its key evaluative terms in ways that make its
central claims tautologically true. A framework exhibits
$\FailGeo{2}$ when it is impossible to specify an
observation that would count as a case of the evaluated
phenomenon occurring without the defining property, because
the defining property is constitutive of what the phenomenon
is, on the framework's own terms. Definitional Closure
renders a framework unfalsifiable not by making its
predictions vague but by making its key terms non-independent
of its conclusions.
\end{auditprop}

\section{The Salience--Meaning Entanglement}
\label{sec:salience-meaning}

The most consequential instance of potential Definitional
Closure in the Phoenix corpus concerns the relationship
between the Marine salience metric and the concept of meaning
or significance. This relationship is central to the corpus's
claim that the Marine detects something genuinely important
in signals rather than just something formally measurable.

The corpus's argument runs roughly as follows: the Marine
detects temporal coherence; temporally coherent signals are
signals that carry meaningful structure; therefore the Marine
detects meaningful structure. The first premise is
operational, the second is the crucial move, and the
conclusion would follow if the second premise were established
independently.

The problem is that the second premise tends to be established
in one of two ways. Sometimes it is established by appeal
to psychoacoustic evidence: coherent signals are easier to
perceive and group, and the auditory system has evolved to
prioritize them. This is an empirical argument, and if the
evidence supports it, it legitimately grounds the connection
between coherence and meaningfulness. Sometimes, however,
the second premise is established by stipulation: meaningful
structure just is the kind of structure that biological
sensory systems prioritize, and biological sensory systems
prioritize temporal coherence, therefore meaningful structure
just is temporal coherence. This is a definitional move,
and it makes the conclusion analytic rather than synthetic.

Whether the corpus is making the empirical or definitional
argument is not always clear, and this ambiguity is itself
a problem. If the argument is empirical, the corpus needs
to provide the evidence in a form that could disconfirm
the connection. If the argument is definitional, the corpus
needs to acknowledge that it has established a terminology
rather than discovered a relationship.

\section{Pseudo-Impossibility Theorems}
\label{sec:pseudo-impossibility}

A particularly interesting manifestation of Definitional
Closure appears in the form of what might be called
\textit{pseudo-impossibility theorems}: claims that have
the logical form of impossibility results but derive their
impossibility from definitional choices rather than from
substantive constraints on the world.

The Phoenix corpus contains several instances. Consider
the claim that conventional audio codecs cannot preserve
the experiential richness of the original signal. This is
presented as an architectural fact about lossy compression:
by removing frequency components, codecs necessarily
destroy something that matters for experience. If this
claim were established empirically --- by demonstrating
that listeners consistently distinguish compressed from
uncompressed audio under conditions where only temporal
structure differs, not overall loudness, spectral centroid,
or other co-varying factors --- it would be a genuine
result.

But the claim also admits a definitional reading: if
``experiential richness'' is defined as ``what temporally
coherent signals provide,'' and if lossy compression
reduces temporal coherence, then it follows by definition
that codecs reduce experiential richness. The impossibility
--- the codec cannot preserve richness --- is a consequence
of the definition rather than a discovery about codecs.

\begin{dangerzone}[Ontological Escalation]
Pseudo-impossibility theorems are one of the mechanisms
by which ontological escalation occurs. A claim that
begins as a specific engineering observation (codecs
disrupt timing relationships) acquires impossibility
status (codecs cannot preserve experience) through a
definitional manoeuvre (experience is defined in terms
of timing relationships). The escalation from specific
to general, from empirical to definitional, happens
smoothly because the terminology persists while its
referent shifts. The result is an apparently strong result
--- an impossibility theorem --- that is actually weaker
than the original engineering observation, because the
impossibility derives from a definition that has not been
justified independently.
\end{dangerzone}

\section{The Constructive Role of Definitions}
\label{sec:constructive-definitions}

Before proceeding, it is important to acknowledge what
Definitional Closure is not. Definitions are essential
to theoretical work; every framework requires definitional
choices, and those choices shape what the framework can
say. The problem is not that the Phoenix corpus makes
definitional choices. The problem is a specific structural
arrangement in which the framework's central evaluative
claims derive their apparent force from definitional
choices rather than from evidence.

Good definitions do three things well. They make a term
precise in a way that is non-circular --- the definition
does not use the term being defined or its synonyms.
They are independently motivated --- there is a reason
to adopt this definition rather than a competitor that
does not depend on the conclusion the definition enables.
And they generate non-trivial consequences --- the
definitional choice commits the framework to claims that
go beyond the definition and that could in principle fail.

By these criteria, the Marine jitter metric is well-defined:
it is precise, operationalizable, and generates non-trivial
predictions about signal behaviour. The term ``authenticity''
as used in the corpus is less well-defined: it is not
clearly operationalized, its motivation appears to depend
on the conclusion it enables (temporal coherence is what
makes audio authentic), and it is not clear what claims
it generates beyond those already made by the jitter metric.

\begin{epistemicstatus}{Reconstructive Conjecture}
The claim that the Phoenix corpus exhibits Definitional
Closure in the strong sense --- that its central claims
are definitionally rather than empirically guaranteed ---
is a conjecture rather than a demonstrated result at this
stage of the analysis. What has been established is a
pattern: several of the corpus's key evaluative terms
(authenticity, meaningfulness, richness) tend to function
in ways consistent with being defined as temporal coherence
under different names. Whether this pattern reflects
genuine Definitional Closure or merely an as-yet-incomplete
attempt to operationalize concepts that have independent
motivation requires the more systematic analysis of Part~III.
\end{epistemicstatus}

\section{The Compression--Cognition Identity}
\label{sec:compression-cognition-identity}

The most philosophically significant instance of potential
Definitional Closure in the corpus concerns the relationship
between compression and cognition. The corpus argues that
cognition is, at its most fundamental level, a compression
operation: it takes high-dimensional, continuous sensory
input and produces a low-dimensional, sparse, temporally
structured representation. This compression is not a
secondary feature of cognition but its essence.

This claim has an empirical reading and a definitional reading.
The empirical reading: cognition, as we observe it in
biological systems and as we wish to implement it in
artificial ones, turns out to involve operations that
can be characterized as compression in the information-theoretic
sense, and understanding this compression-theoretic character
helps us understand how cognition works. This is a
legitimate scientific hypothesis that generates testable
predictions.

The definitional reading: cognition just is a certain form
of compression; anything that performs this compression
is cognitive; anything cognitive performs this compression.
On this reading, the claim that ``the Marine detects
meaningful structure because meaningful structure is
what efficient compression targets'' is analytic rather
than synthetic. It cannot be falsified because it is
built into the definition of what meaningful structure is.

The corpus appears to move between these two readings
without clearly committing to either. This ambiguity is
not a deliberate evasion; it reflects the genuine
difficulty of distinguishing, at an early stage of
framework development, between a framework that has chosen
a productive definition and a framework that has built its
conclusions into its definitions. But the ambiguity has
epistemic consequences: it makes it impossible, from
outside the framework, to determine whether the corpus
is making discoveries or stipulations.

\section{What Closure Would Prevent}
\label{sec:closure-prevents}

Definitional Closure, if present, would prevent the Phoenix
corpus from doing something that it clearly wants to do:
engage with evidence. The corpus repeatedly refers to
neuroscientific findings, psychoacoustic research, and
engineering benchmarks as sources of support for its claims.
This referential behaviour implies that the corpus intends
its claims to be empirical rather than definitional ---
that it intends them to be claims that the evidence could
in principle disconfirm.

This is worth emphasizing. The corpus is not claiming to
be providing definitions; it is claiming to be making
discoveries. If those claims are actually definitional,
the corpus has not achieved what it set out to achieve.
The fault is not malice but structure: a framework that
builds its conclusions into its definitional choices may
do so without recognizing that this is what it has done,
and the recognition of what has happened requires looking
at the framework from a position outside it.

This is one of the most important functions of the \Ortyx{}
Protocol as a meta-methodology: to provide a vantage point
outside any given framework from which to ask whether
the framework's claims are doing the work they appear to
be doing. Whether Ortyx itself maintains this vantage point
--- whether it avoids building its own conclusions into
its audit definitions --- is a question for Part~V.

\medskip

Chapter~9 examines the third failure geometry: Semantic
Universality Collapse, which occurs when a framework's
abstraction level rises to the point where its central
concepts can accommodate any observation. This is
the dynamic that produces what Part~I called the
``universal absorbency'' risk: the framework that
explains everything, forbids nothing, and thereby ceases
to explain anything at all. The Phoenix corpus's expansive
deployment of temporal coherence as a universal organizing
principle makes it particularly susceptible to this failure
mode.
